% !TEX root = ../main.tex

\section{Related Work}
\label{ch:related-work}

% Other approaches to the problem considered reported in the literature. Comparison to the paper at hand.
\textbf{Post-Hoc Methods:} These are used to analyze and interpret machine learning models after they have been trained. \textcite{clark2019does} analyses the attention patterns and \textcite{tenney2019bert,belinkov2022probing} tried to probe the hidden representaions i.e. train a classifier on the hidden representations to predict the properties learned by the model. But it does not tell if the model use the property discovered by probe for prediction or not. Our work is different as our model is interpretable by design and does not require any post-hoc analysis.

\textbf{Program Induction and Neuro-Symbolic Methods:} These approaches focus on learning rule-based programs from data, often leveraging neural networks to guide the search for discrete logical expressions, as demonstrated by \textcite{wang2021scalable,petersen2022deep}. While both our work and these methods create interpretable discrete models, the key difference lies in their scope—these methods do not target Transformers or sequence modeling tasks, which are central to this work.  


\textbf{Transformers and formal languages:} This area of research try to focus on theoretical capabilities of transformer with hard attention such as their expressive power \cite{hao2022formal, merrill2022saturated}. For instance, \textcite{merrill2022saturated} demonstrated that Transformers could function as programmable computers by simulating basic calculator and different linear algebra operations. However, this line of work remains theoretical. Furthermore, \textcite{merrill2024logic} proved that any fixed-precision Transformer could be expressed as a formula within a first-order logic framework. In contrast to these studies, our work emphasizes developing a Transformer architecture that prioritizes both interpretability and practical application, rather than extracting rules from pre-trained models.

\textbf{Designing Interpretable Models:} Several studies have proposed modifications to Transformer architectures to improve their interpretability. For example, sparse attention mechanisms by \textcite{zhang2021sparse} and adjustments to activation functions by \textcite{elhage2022solu} have been introduced to make specific components easier to analyze. However, these approaches focus on enhancing the qualitative understanding of certain elements rather than creating fully interpretable models.